\section{Actual exposure by cofactor--orbit cells}

The row-singleton theorem identifies outputs with literal atoms, but it does
not prove that a row contains any such atom. We record the exact exposure
decomposition and its coefficient-blind upper bound.

For an active row \(m\), a cofactor \(r\in[R,2R]\), and an orbit time \(j\),
put
\[
 \Gamma_m(r,j)
 :=\{\gamma:(m,r,j,\gamma)\in\mathcal U_X(R)\}
\]
and
\begin{equation}\label{eq:tpc66-cell-energy}
 E_m(r,j):=\sum_{\gamma\in\Gamma_m(r,j)}
 |a_{m,r,j,\gamma}|^2.
\end{equation}
Then the exposure and energy disintegrate exactly:
\begin{align}
 S_m
 &=\bigsqcup_{\substack{r\in[R,2R]\\j}}
   \{(\gamma,d_mr,j):\gamma\in\Gamma_m(r,j)\},
 \label{eq:tpc66-exposure-disintegration}\\
 T_m(R)&=\sum_{r,j}E_m(r,j).
 \label{eq:tpc66-energy-disintegration}
\end{align}

\begin{proposition}[One-row orbit spacing]\label{prop:tpc66-row-spacing}
For sufficiently large \(X\), fix a literal row \(m=d_m\ell_m\) and a
retained cofactor \(r\). If one solution of
\[
 rp-mj=h_0
\]
exists, all integral solutions are
\begin{equation}\label{eq:tpc66-row-ladder}
 j=j_0+rt,\qquad p=p_0+mt,\qquad t\in\mathbb Z.
\end{equation}
Consequently an orbit interval of length \(O(J)\) contains at most
\begin{equation}\label{eq:tpc66-row-exposure-upper}
 1+O\!\left(\frac Jr\right)
 \le1+O\!\left(\frac JR\right)
\end{equation}
candidate \(j\)-values for this \((m,r)\).
\end{proposition}

\begin{proof}
On the inherited terminal support, \((d_m,r)=1\). If the source prime
\(\ell_m\) divided \(r\), then it would divide both \(rp\) and
\(d_m\ell_mj\), hence \(h_0\), impossible once
\(\ell_m>|h_0|\). Thus \((m,r)=1\). Differences of two solutions obey
\(r\Delta p=m\Delta j\), yielding \eqref{eq:tpc66-row-ladder}. The spacing
of \(j\) is \(r\), which gives \eqref{eq:tpc66-row-exposure-upper}.
\end{proof}

Primality of \(p\), the mask, the Mellin profile, and every other literal
support condition only delete these candidates. Therefore
\eqref{eq:tpc66-row-exposure-upper} is an upper census. It proves neither
nonemptiness nor a lower bound for \(T_m(R)\). Even a large candidate count
does not imply energy diffuseness among the surviving cells.
